
\pagebreak

\chapter{Web performance optimization, lazy loading.}


\medskip

Web performance optimization is a crucial aspect of modern development because the speed and responsiveness of a website directly influence how users perceive it and whether they choose to stay or leave. One of the most effective techniques for improving performance is lazy loading. The idea behind lazy loading is simple: instead of forcing the browser to load every single asset or piece of content immediately when a page opens, resources are loaded only when they are needed. This reduces the initial loading time and improves the overall user experience.

Consider a webpage that contains dozens of high-resolution images. If all of them are downloaded at once, the page will take longer to render, especially on slower networks or mobile devices. This can frustrate users who might leave before even seeing the content. With lazy loading, only the images visible in the user�s viewport are fetched, while others are loaded later as the user scrolls down. This means the page feels faster and more responsive right from the beginning.

Lazy loading is also beneficial for conserving bandwidth and server resources. For example, on a news site, a user may only read the first few articles and never scroll further. Without lazy loading, all the images and media for the entire page are downloaded unnecessarily, wasting both client and server capacity. By deferring these requests, the system becomes more efficient and scales better under heavy traffic.

In practice, lazy loading has become particularly important in the context of mobile browsing, where network conditions are often less reliable. A fast and responsive mobile experience can make the difference between keeping a user engaged or losing them to a competitor�s site. That is why lazy loading is not just a technical trick but a user-centered strategy that aligns with the broader goal of web performance optimization.

\medskip

\textbf{Example.}

Imagine you have an endpoint that returns product data including images. Without optimization, you might return everything at once: product details, specifications, and all images in full resolution. This is wasteful, especially if the user only needs the list of products to begin with.

A lazy-loading approach would be to create two endpoints. The first one delivers lightweight product data such as name, price, and a small thumbnail. The second endpoint provides detailed product information, including high-resolution images, but it is only called when the user requests to view the product details.

\begin{lstlisting}
// Lightweight endpoint
[HttpGet("products")]
public IActionResult GetProducts()
{
    var products = _context.Products
        .Select(p => new {
            p.Id,
            p.Name,
            p.Price,
            ThumbnailUrl = p.Images.FirstOrDefault(i => i.IsThumbnail).Url
        })
        .ToList();

    return Ok(products);
}

// Detailed endpoint
[HttpGet("products/{id}")]
public IActionResult GetProductDetails(int id)
{
    var product = _context.Products
        .Where(p => p.Id == id)
        .Select(p => new {
            p.Id,
            p.Name,
            p.Price,
            p.Description,
            Images = p.Images.Select(i => i.Url)
        })
        .FirstOrDefault();

    return Ok(product);
}

\end{lstlisting}

On the frontend, lazy loading ensures that large assets like images are loaded only when they become visible. Modern browsers support this directly with the loading="lazy" attribute:

\begin{lstlisting}
<img src="product-thumbnail.jpg"
     data-full="product-large.jpg"
     loading="lazy"
     alt="Product Image">

\end{lstlisting}

Or with JavaScript using an Intersection Observer to load images only when they enter the viewport:

\begin{lstlisting}
const images = document.querySelectorAll('img[data-src]');

const lazyLoad = (entry, observer) => {
  entry.forEach(e => {
    if (e.isIntersecting) {
      const img = e.target;
      img.src = img.dataset.src;
      observer.unobserve(img);
    }
  });
};

const observer = new IntersectionObserver(lazyLoad, { threshold: 0.1 });

images.forEach(img => observer.observe(img));

\end{lstlisting}

